(1) 对传输线微元作电压降与电流连续性平衡，得 $V_x=-L'I_t$、$I_x=-C'V_t$。分别再求空间导数并消去另一变量，便有 $V_{xx}=L'C'V_{tt}$ 和 $I_{xx}=L'C'I_{tt}$，说明波形在理想无耗线路上无衰减地传播。
(2) 由标准一维波动方程可知 $v=1/\sqrt{L'C'}$。沿 $+x$ 的电压行波 $V^+=f(x-vt)$ 满足 $I^+=V^+/Z_0$，其中 $Z_0=\sqrt{L'/C'}$；沿 $-x$ 的电压行波则满足 $I^-=-V^-/Z_0$，负号反映能流方向反转。
(3) 负载端电压为 $V_L=V_i+V_r$，电流为 $I_L=(V_i-V_r)/Z_0$。以 $V_L=R_LI_L$ 联立解得 $\Gamma_V=V_r/V_i=(R_L-Z_0)/(R_L+Z_0)$，而电流反射系数 $\Gamma_I=-\Gamma_V$，符合右行和左行波电流方向相反的约定。
(4) 无损传输线平均功率与电压峰值平方成正比，因此反射功率比 $|\Gamma_V|^2$，负载吸收比 $1-|\Gamma_V|^2=4R_LZ_0/(R_L+Z_0)^2$。匹配 $R_L=Z_0$ 时零反射；开路给 $\Gamma_V=+1$，短路给 $\Gamma_V=-1$，二者均不吸收净功率。
